\begin{frame}
    \frametitle{Le reti nel mondo reale}
    Molti sistemi reali possono essere rappresentati con una rete.
    \begin{columns}
        \begin{column}{0.4\textwidth}
            \begin{exampleblock}{Alcuni esempi}
                \begin{enumerate}
                    \item Reti sociali
                    \item Reti dei trasporti
                    \item Reti biologiche
                \end{enumerate}
            \end{exampleblock}
            % \pause
            \begin{alertblock}{Problema}
                Un singolo \alert{grafo} permette di rappresentare un unico tipo di relazione tra le entità
            \end{alertblock}
        \end{column}
        \begin{column}{0.4\textwidth}
            \begin{figure}
                \centering
                \includegraphics[width=\textwidth]{img/graph.pdf}
                % \caption{Grafo aggregato derivato dalla multilayer network rappresentata in Figura \ref{fig:mlexample}.}
                % \label{fig:graph}
            \end{figure}
        \end{column}    
    \end{columns}
   
    % \begin{definizione}Un 
    % \alert{numero primo} è un intero $>1$ che ha esattamente due divisori positivi.
    % \end{definizione}
\end{frame}

% \begin{frame}
%     \frametitle{I numeri primi sono infiniti}
%     \framesubtitle{Ne diamo una dimostrazione diretta}
%     \begin{teorema}Non esiste un primo maggiore di tutti gli altri.
% \end{teorema}
%     \pause
%     \begin{proof}
%     \begin{enumerate}[<+->]
%     \item Sia dato un elenco di primi.
%     \item Sia $q$ il loro prodotto.
%     \item Allora $q+1$ è divisibile per un primo $p$ che non compare nell'elenco. 
%     \qedhere
% \end{enumerate}
% \end{proof}
% \end{frame}
% \begin{frame}
%     \frametitle{Che cosa c'è ancora da fare?}
%     \begin{block}{Problemi risolti}Quanti sono i numeri primi?
% \end{block}
%     \begin{block}{Problemi aperti}Un numero pari $>2$ è sempre la somma di due primi?
% \end{block}
% \end{frame}